\section{Neural Field：同一个 Whole 如何生成不同 Parts}

最后两分钟补上了一个容易被遗漏、却决定 top-down path 是否自洽的问题：若 face-level island 在所有位置都是同一 vector，而且所有 column 共享同一个 top-down network，为什么有的位置输出 nose、有的位置输出 mouth？答案是 top-down function 还接收 target location。

\subsection{Location-conditioned top-down prediction}

令 whole embedding 为 $\mathbf e_w$，需要预测的位置为 $x$，则 shared decoder 是函数

\[
\hat{\mathbf e}_{\mathrm{part}}(x)=f_{\mathrm{td}}\!\left(\mathbf e_w,x\right).
\]

其中，$\mathbf e_w$ 同时编码 whole identity 与 pose，$x$ 指定要查询的图像位置，输出是在该位置应出现的 lower-level identity-pose embedding。相同 face vector 配合不同 $x$，可以生成 nose、mouth 或其他 part prediction；这就是 coordinate-conditioned neural field 的基本形式。

\teachervoice{Hinton 先提出一个表面矛盾：red arrows 与 green arrows 完全不同，上层 face vectors 却相同；如果 top-down weights 也共享，怎么产生不同结果？他的解答是把 query location 一并输入。这个补充把“共享权重”从口号变成了函数签名。}

\begin{knowledgebox}{Neural field 在这里解决什么}
\begin{itemize}
  \item 输入：whole identity/pose embedding 加 target coordinate；
  \item 输出：该 coordinate 上的 lower-level prediction；
  \item 价值：共享同一 decoder，却能在不同位置产生不同部件；
  \item 不等价：课堂并未要求使用 NeRF 的 volume rendering、ray sampling 或 photometric loss。
\end{itemize}
\end{knowledgebox}

\subsection{与现代 coordinate-conditioned model 的连接}

NeRF、implicit neural representation、一些 world model decoder 都使用“global/local latent + coordinate/query”产生位置特定输出。GLOM 的独特点在于 decoder 不是只从一个 global latent 渲染像素，而嵌在多层 recurrent islands 中：top-down prediction 必须与 bottom-up evidence、lateral agreement 和 temporal state 共同收敛。

\begin{warningbox}{Connection 不是 lineage claim}
把 GLOM 与现代 object-centric decoder、slot model、world model 相连，是工程解释；不能说后者“实现了 GLOM”，除非它们真的满足 island representation、multi-level recurrence、four-source consensus 与 coordinate-aware part-whole prediction 等更强条件。
\end{warningbox}

\subsection{本章小结}

Location condition 解决了 shared top-down network 的表达矛盾：同一个 whole representation 可以针对不同坐标发出不同 part prediction。它也是整套设计中最接近成熟现代技术的组件，但 GLOM 的难点仍在于如何把该 decoder 与动态 islands 稳定联合训练。

